%!TEX root=../a03.tex
\section{Q2}
Prove the following statement. 
\begin{equation*}
    \forall x \in \N, \exists y \in \N, \paren{x^2 + 2 - y} \mid \paren{10 - x^2 - y^2}
\end{equation*}
\begin{solution}
    Let $x$ be arbitrary number in $\N$. We select $y$ to be $3$ in $\N$. 

    subtituting $y = 3$ into the expression, we 
    \begin{equation*}
        \paren{x^2 + 2 - 3} \mid \paren{10 - x^2 - \paren{3}^2}
    \end{equation*}

    This implies 
    \begin{equation*}
        \paren{x^2 - 1} \mid \paren{10 - x^2 - 9 }, 
    \end{equation*}

    which equals to 
    \begin{equation*}
        \paren{x^2 - 1} \mid \paren{1 - x^2}
    \end{equation*}

    there exists a number $k = -1$ in $\Z$ such that 
    \begin{equation*}
        k(x^2 - 1) = -1\paren{x^2 - 1} = -x^2 + 1 = 1 - x^2. 
    \end{equation*}

    for all $x \in \N$. 

    Thus, according to the definition of divisibility, 
    \begin{equation*}
        \forall x \in \N, \exists y \in \N, \paren{x^2 + 2x - y} \mid \paren{10 - x^2 - y^2}. 
    \end{equation*}
\end{solution}